{
\def\UrlBreaks{\do a\do b\do c\do d\do e\do f\do g\do h\do i\do j\do k\do l\do m\do n\do o\do p\do q\do r\do s\do t\do u\do v\do w\do x\do y\do z\do A\do B\do C\do D\do E\do F\do G\do H\do I\do J\do K\do L\do M\do N\do O\do P\do Q\do R\do S\do T\do U\do V\do W\do X\do Y\do Z\do 0\do 1\do 2\do 3\do 4\do 5\do 6\do 7\do 8\do 9\do /\do _\do .\do -}
\newcommand{\CertificateInterval}[2]{\left[\begin{gathered}{}#1,\\{}#2\end{gathered}\right]}
\section{Numerical certificates}

The source comparisons hold by Sections~2--5.
Section~6 uses each selected inequality once,
with the signs of (6.1).
The tables give full rational certificates and independent checks.
Their directed interval calculations use Arb's midpoint-radius arithmetic
\cite{Johansson2017}; rational endpoint checks are given explicitly.

\subsection{Exact interval and normalization conventions}

Each terminating decimal is an exact rational,
and each endpoint includes its displayed radius.
The extra digits fix the direction of each bound.
Set
\[
 B(c,r)=[c-r,c+r].                                      \tag{A.1}
\]
Shorter intervals are rounded outward.
The numerical-certificate supplement archive \path{paper-ii-public-certificates-v11.tar.gz} has SHA-256
\path{cb9d7b96ce1323c9118d845f4dd701a0d830e8c610090deec88bd8792648e100}.
Each printed file SHA-256 is the digest of that file in the archive.
An endpoint certificate bounds its endpoint,
and a functional enclosure bounds the integral.
Paths are relative to the extracted numerical-certificate supplement.

The common parameters are
\[
\begin{aligned}
 u&=\frac1{500},\quad \alpha=\frac{100}{1149},\quad
 \beta=\frac{50}{309},\\
 \mu&=\frac{25}{128},\quad \lambda=\frac{93}{400},\quad
 \tau_S=\frac{57}{224},\\
 t_*&=\frac{7109}{28000}.
\end{aligned}\tag{A.2}
\]
The parameter-family certificates require
$0<\sigma\le10^{-10}$ and $\tau=\sigma^9\le10^{-90}$.
They apply to the same source and scalar profile.
In Section~4 the G1 source width is $C_b=334$.

The two numerical outputs and their component manifests are:

{
\normalsize
\renewcommand{\arraystretch}{1.14}
\begin{longtable}{@{}p{\textwidth}@{}}
\toprule
\textbf{\path{endpoint-recert/out/compare_hp.json}, S1, row and G1} \\
\emph{SHA-256}: \path{cff4d6dc0168e66a29a915473c3391c383c168fb81ac71ecec87e90016331bf8} \\
\emph{Manifest SHA-256}: \path{8760a5a78fff44ca739cd12660d3a86765dd2f3d6872eae111d933f354f54730} \\
\midrule
\textbf{\path{s2-joint/outputs/ASSEMBLY.json}, S2} \\
\emph{SHA-256}: \path{21083883403f3903fff3ce781e086a26f931dd809e445578cfcba546295e9e5d} \\
\emph{Manifest SHA-256}: \path{d26488e736e5798544e758338acb299d417b74a309568f5ce20f738e399cf96d} \\
\bottomrule
\end{longtable}
}

Certificate paths are relative to the supplement root.
The root \path{SHA256SUMS} lists each retained filename and digest,
including the selected S2 components.
The true $F/f$ profiles are computed by \path{endpoint-recert/code/tm.py}, with SHA-256
\path{d34b11f02353473f714d00dbb02d73f9a91be5275d9589fc0db1de3aa851ec16}.
The exact comparisons are in \path{endpoint-recert/out/compare_hp.json}, with SHA-256
\path{cff4d6dc0168e66a29a915473c3391c383c168fb81ac71ecec87e90016331bf8}.
The bounds come from Taylor models and directed quadrature.

Throughout, $U_C=C(N)N/\log^2N$, with $C(N)$ as in \S~1.
The signed master coefficients, in the order
$S_1,\ldots,S_{13},S_{\rm sw}$, are
\[
 (3,1,-2,-1,-1,-1,-1,1,1,1,-2,-1,-1,-1).               \tag{A.3}
\]
Thus S1 contributes $3S_1/4$ and S3 contributes $-S_3/2$.
The multiplier is $-1/4$ for S4/S5/S6/S7
and $1/4$ for S8/S9/S10.
The joint S2 and G1 values are already in D units.
The G1 source multiplicities $(2,1,1)$ sit inside its source weight,
so they add no DIRECT-row coefficients.
The switched term contributes $-S_{\rm sw}/\mathcal W$.
Up to its stated vanishing error, this is bounded below by $-2H_L$.

Each selected endpoint has an independent check,
and each row and joint contribution enters once.

\subsection{S1: joint enclosure}

For the joint profile in Section~2,
\[
 12.9020516<S_1<12.9023390,\qquad
 \frac34S_1>9.6765387.                                 \tag{A.4}
\]
It covers S1-0 through S1-3 and includes the first $-F$ subtraction.
The subtraction uses $u$ and $s_0=19037/64$.
The integral of Section~2.3 satisfies
\[ I_*<2161.074150710003. \]
Both bounds refer to the joint profile of Section~2.3.

The selected enclosure is computed by
\path{s1-tail/certify_s1_corrected_tail.py},
with SHA-256
\path{da9f572cbb28ca06a00e050785eefd9cd4f499269a627deaaa9a593c4ab210a3};
its output
\path{s1-tail/corrected-certificate.json}
has SHA-256
\path{ebf20d49c73cbf7ae11314e3a1e7e82005a75c42a6c2fe91547698f877011f4f}.
The infinite factorial tail is
\[
 \sum_{m\ge40}\frac{(m+2)^2}{(m+1)m!}
 \le\frac{42^2}{41\,40!}\frac{39}{38},                 \tag{A.5}
\]
since the ratio of consecutive summands is at most $1/39$.
The endpoint (A.4) is used.

{
\normalsize
\renewcommand{\arraystretch}{1.14}
\begin{longtable}{@{}p{\textwidth}@{}}
\toprule
\emph{Exact outer intervals containing the computed bounds} \\
\midrule
\textbf{S1 lower bound} \\
$B(12.9020516652549921721098499909159212126626402875924361285,4.64\cdot10^{-56})$ \\
\midrule
\textbf{S1 upper bound} \\
$B(12.9023389323670212025588110758944404798760246915400198042,3.39\cdot10^{-56})$ \\
\midrule
\textbf{D lower bound from S1} \\
$B(9.6765387489412441290823874931869409094969802156943270964,6.00\cdot10^{-56})$ \\
\midrule
\textbf{Infinite-tail majorant $J_{\rm upper}$} \\
$B(0.000255821764915303175726784912190615770902363950678941419476,3.34\cdot10^{-61})$ \\
\midrule
\textbf{Prefix deficit} \\
$[0,7.792\cdot10^{-383}]$ \\
\bottomrule
\end{longtable}
}

The prefix deficit is subtracted before computing the lower endpoints in (A.4).
Thus (A.4) includes the deficit once.
The source-to-profile comparison is in \S~2.3.

The independent calculation is \path{endpoint-recert/code/s1_cert.py}, with SHA-256
\path{718e495fb18ea71ef0538c73ac03f96a0cac77f2274ffbe5c1a093b76605a508}.
It re-derives the cancellation and tail bound from the delay equations.
Its output is \path{endpoint-recert/out/s1_cert.json}, with SHA-256
\path{f16972e2b73258532182357454193b5002b553dc465b829f9af9b2c0680d4137}.
It gives the sharper bounds
\[
 S_1\in[12.9020760,12.9023390],\qquad
 \frac34S_1>9.6765570.                                 \tag{A.5a}
\]
The computed lower-bound ball is
{\normalsize
\[
B(12.9020760472276414637285015429505331626940241680712780695,
4.38\cdot10^{-56})\mathpunct{,}
\]
and its $3/4$ multiple is
\[
B(9.67655703542073109779637615721289987202051812605345855210,
5.80\cdot10^{-56})\mathpunct{.}
\]
}
These bounds imply (A.4) and its tail estimate.
They also give $I_*<2161.074150710003$.
The selected endpoint is (A.4).

\subsection{The seven separate profile endpoints}

The seven directed endpoints below come from three certificate groups.
Option-A gives the absolute S3-A/S8-A/S9-A values.
Direct-row gives S4, S5-low and S7,
and Retained-row gives S10-Ret.
Each endpoint has an independent comparison.
For S4, S5-low, S7 and S10 the final sum uses the adverse endpoint of the
first printed interval; the following ball is an independent check.
For the three Option-A rows it uses the displayed ball endpoint in the
adverse direction.
S7 is taken at $\lambda=93/400$.
S10 uses the source-trimmed upper inner boundary.

{
\normalsize
\renewcommand{\arraystretch}{1.14}
\begin{longtable}{@{}p{\textwidth}@{}}
\toprule
\emph{Row; coefficient giving its D contribution; certificate group.} \\
\emph{Raw-profile outer interval; next line: independent endpoint comparison.} \\
\midrule
\textbf{S3-A (Option A)}\quad $-1/2$\quad Option-A \\
$B(10.45234847108890216860085651838875393044983358511004,5.69\cdot10^{-51})$ \\
Upper endpoint verified in the archived comparison record \\
\midrule
\textbf{S4}\quad $-1/4$\quad Direct-row \\
$[3.336569865049636,3.336975348289866]$ \\
$B(3.33672680275562584206999257242,4.25\cdot10^{-30})$ \\
\midrule
\textbf{S5-low}\quad $-1/4$\quad Direct-row, direct cells $[1/4,t_*)$ \\
$[0.238854512468062,0.238888433191047]$ \\
$B(0.238875420270830663069285634601,4.07\cdot10^{-31})$ \\
\midrule
\textbf{S7}\quad $-1/4$\quad Direct-row, $\lambda=93/400$ \\
$[2.299209653294856,2.299402520476652]$ \\
$B(2.29926042756192568371844970138,1.46\cdot10^{-30})$ \\
\midrule
\textbf{S8-A (Option A)}\quad $+1/4$\quad Option-A \\
$B(2.36595898835331722398869549955202092760264561327680,6.59\cdot10^{-51})$ \\
Lower endpoint verified in the archived comparison record \\
\midrule
\textbf{S9-A (Option A)}\quad $+1/4$\quad Option-A \\
$B(1.341531658234294936348198541015408702983637197761145323037489950,4.21\cdot10^{-64})$ \\
Lower endpoint verified in the archived comparison record \\
\midrule
\textbf{S10-Ret}\quad $+1/4$\quad Retained-row \\
$[0.858847455777940,0.867486588163666]$ \\
$B(0.862377086132397387439154957615,6.98\cdot10^{-32})$ \\
\bottomrule
\end{longtable}
}

{
\normalsize
\renewcommand{\arraystretch}{1.14}
\begin{longtable}{@{}p{\textwidth}@{}}
\toprule
\textbf{Direct-row} \\
\emph{Source code and SHA-256}: \path{row-enclosures/cert_s4_s7.py}: \path{d98866b45f1ddb8f3f93411bfe09cb11c62569cc80fa775a77a3eab1ad4384aa} \\
\emph{Output and SHA-256}: \path{row-enclosures/result.json}: \path{0eb5fa52ed180a071633e50f0d099f5c4e1f9bb001add71765fe1c2f033b8f7a}; S4, S5 direct cells, S7 \\
\midrule
\textbf{Option-A} \\
\emph{Source code and SHA-256}: \path{option-a/arb_enclose_optionA_independent.py}: \path{9c1880609692d30be9f282a218280c291dbb782a7d7cfab35732e9203687c36b} \\
\emph{Output and SHA-256}: \path{option-a/arb_enclose_optionA_out.json}: \path{5a8786a27da1a52550eb0c077ae11c48072aff6d1765cebbad25b24e2a556b27} \\
\midrule
\textbf{Retained-row} \\
\emph{Source code and SHA-256}: \path{s10/arb_s10_exact.py}: \path{74a130143e5c9e40ac530dd688526e63acde93a95f5e8271b0f07984afc5a98f} \\
\emph{Output and SHA-256}: \path{s10/s10-s10.json}: \path{f6fba094d222f362341bc63ba0388cbf6a7556accdd14d5b28b882d7e0626cf6} \\
\bottomrule
\end{longtable}
}

Direct-row uses 192-bit Arb with $\tau\in[0,10^{-90}]$.
S5-low sums two direct cells on $[1/4,t_*)$, with outward rounding.
The BV-top cell is disjoint from them. The independent
calculation
\path{endpoint-recert/code/rows_cert.py} has SHA-256
\path{8b2ef756bbdd5530f614a66b8c5eb5dddd17922555d979515b1ca85876064bc5}.
It produces all seven results in \path{endpoint-recert/out/rows_cert_hp.json}, with SHA-256
\path{456d03d3ea9954bb71fedc8e20ea35cd463454bb9a0919b3ace9f64e09f54464}.
The exact rational comparisons are in \path{endpoint-recert/out/compare_hp.json} of Appendix A.1.
They include the selected S8/S9 directions.
S3/S4/S5/S7 include the adverse $\tau=10^{-90}$ effects,
and S8/S9/S10 use the actual-coordinate profiles.
Their fixed-$\tau$ comparisons are proved in Sections~2.4--2.5,
especially (2.40).

The independent S10 retained-strip bound is $4.82\tau$ in D units.
It lies within the single separate $7.160\tau$ bound.

The absolute Option-A rows use one source,
with the subtraction at $u=1/500$.
These rows give the three selected absolute endpoints.

The independent comparison gives
{\normalsize
\[
\begin{gathered}
 S_{3,\mathrm{OptA}}\in\\
B(10.45234847108890216860085651838875393044983358511004,
   5.69\cdot10^{-51}),
\end{gathered}
\]
}
and the matched change is
{\normalsize
\[
\begin{gathered}
 \Delta D_{\mathrm{OptA}}\in\\
B(-0.0110499568842910017666362322996150028835154631047,
   9.53\cdot10^{-50}).
\end{gathered}\tag{A.6}
\]
}
Equation (A.6) is a comparison value only.
The signed sum uses the three absolute endpoints.

\subsection{S2: one four-row endpoint}

Let $P_-$, $J_+$ and $C_-$ be directed endpoints.
They bound the prefix below, PF above and the BV corner below.
On $\Delta=\{u\le r_3<r_2<r_1\le\beta\}$, set
\[
 V_{\rm full}=V
 =\int_\Delta\frac{dr_1\,dr_2\,dr_3}{r_1r_2r_3^2}
 =\frac{\log^2(\beta/u)-2\log(\beta/u)+2}{2u}-\frac1\beta.
\]
An outward evaluation gives
{\small
\[
\begin{aligned}
 V_{\rm full}\in[&
3122.423962424988515509469383070552351626538598751986375091552742,\\
&3122.423962424988515509469383070552351626538598751986375091552743].
\end{aligned}
\]
}
The endpoints in (A.7) are
\[
\begin{aligned}
P_-&=3513.1687300568893335013223065223679611140065161989243,\\
J_+&=0.359440819140453278302418897934329259372468992946659201284,\\
C_-&=-0.0574908755554524530180814901226778275854309897822443970870.
\end{aligned}
\]
Use four times the printed interval for $1/(2e^\gamma)$ below to bound
$2/e^\gamma$. Define $d_{T8}$ by
\[
 d_{T8}=\frac{P_--(2/e^\gamma)V_{\rm full}
                    -(2/e^\gamma)J_++C_-}{4}
              -\frac1{20000}.                         \tag{A.7}
\]
The directed assembly gives
\[
 d_{T8}>
 1.61958135698908352559133404332105729220076714081650.    \tag{A.8}
\]
Let $k_-,k_+$ be the printed endpoints for $k=1/(2e^\gamma)$ below,
and let $V_-,V_+$ be the endpoints above. Exact rational interval
arithmetic in (A.7) gives
\[
\begin{aligned}
 d_{T8}\in\bigl[&P_-/4-k_+(V_++J_+)+C_-/4-1/20000,\\
                 &P_-/4-k_-(V_-+J_+)+C_-/4-1/20000\bigr].
\end{aligned}
\]
This is the selected joint lower value in D units.
It covers S2-0 through S2-3(K/PF)
and includes the tail $1/20000$ once.
The term $275009680\tau$ is separate,
and so are AA37 and the other finite-$N$ errors.
The grid is the global one, $\tau_g=\tau=\sigma^9$.
The selector and PR costs satisfy
\[
\begin{aligned}
 &(275009680+7.160+62500+145664566.4567)\tau\\
 &=420736753.6167\tau\le4.207367536167\cdot10^{-82}.
\end{aligned}
\]
The BV corner is $c_\gamma V_B-\mathrm{IBV}$ from (5.18). IBV is
the R1 kernel (5.17) at $\kappa=10^{-6}$ in \S~5.3. The independent
nested Arb source is \path{s2-joint/code/corner.py}, with SHA-256
\path{7cbc5d276b54b52d657fc52622db1a3542c5079d72a39e3be3d12eec1ddb21f2}.
Its line 11 sets $\kappa=10^{-6}$. It gives
\[
 c_\gamma V_B-\mathrm{IBV}\in
 B(-0.057257305790446009369486009082468758,4.67\cdot10^{-37}).
\]
The output is \path{s2-joint/outputs/corner.json}, with SHA-256
\path{360841b6008cb08be571d23177887df14be286ebde7ba3f3fe2c7571f31be9d8}.
This value exceeds $C_-$ by more than $2.3356\cdot10^{-4}$, so $C_-$
is a lower endpoint at $\kappa=10^{-6}$. Increasing $\kappa$
decreases $1/2-\kappa-r_1-r_2-r_3$ and lowers the corner, so its
$\kappa=0$ value is larger.
For the R1 level increment, the S2 corner uses $\kappa=10^{-6}$,
so $d_{T8}$ includes this increment.
Section~\ref{sec:proof-theorem-1} keeps an additional conservative R1 allowance.

\paragraph{Corner increment bound.} R1 is the tetrahedron
$\{r_3\le r_2\le r_1\le\beta,\ r_1+2r_2+4r_3\ge H\}$. Its vertices are
$(\beta,\beta,\beta)$, $(H/7,H/7,H/7)$, $(\beta,x,x)$, and
$(\beta,\beta,y)$, where $x=(H-\beta)/6$ and
$y=(H-3\beta)/4=13147/103000$. On R1, $r_i\ge y$ and
$a_0:=1/2-r_1-r_2-r_3\ge3/206$. In D units,
\[
\begin{aligned}
 J_{R1}(\kappa)-J_{R1}(0)&=\int_{R1}\frac{\kappa\,dr}{r_1r_2r_3\,a_0(a_0-\kappa)}\\
&\le\frac{\mathrm{vol}(R1)\,\kappa}{y^3\,\tfrac3{206}(\tfrac3{206}-\kappa)}\\
 &=5.74468\ldots\cdot10^{-6}<5.745\cdot10^{-6}
\end{aligned}
\]
For $\kappa=10^{-6}$,
$\mathrm{vol}(R1)=2.53345\ldots\cdot10^{-6}$. Its exact rational value
is the absolute determinant of three edge vectors at a vertex divided by $6$.
This bounds the actual R1 increment, which is already included in $d_{T8}$.
The extra allowance $g_{R1}=5.71\cdot10^{-6}$ in Section~6 is a separate
conservative subtraction; it is not used as a bound for this increment.

{
\normalsize
\renewcommand{\arraystretch}{1.14}
\begin{longtable}{@{}p{\textwidth}@{}}
\toprule
\emph{Exact outward enclosures} \\
\midrule
\textbf{Prefix, n=600} \\
$P\in\CertificateInterval{3513.1687300568893335013223065223679611140065161989243}{3513.1687300568893335013223065223679611140065161989248}$ \\
\emph{Source code SHA-256}: \path{s2-selected/sources/numeric/split_prefix_exact.py}: \path{492ef149b0460f59732102ddeb252fa8fb30fce03df2fd953376ad4fb5417ada} \\
\emph{Output SHA-256}: \path{s2-selected/sources/numeric/prefix-u-n600-corrected.json}: \path{01c72d03c23ff51eb903e299ecb63a90496485c0e6549e05c3818077f2e88e08} \\
\midrule
\textbf{PF outer, n=1536} \\*
$J\in\CertificateInterval{0.352576588183432872639904522863463727436871987468524955851}{0.359440819140453278302418897934329259372468992946659201284}$ \\
\emph{Source code SHA-256}: \path{s2-selected/sources/numeric/exact_outer_s2.py}: \path{1cf453f317366d9c0b191aac5457b518fae48b54d7321da80bf577d1aa5f69f0}; \path{s2-selected/sources/numeric/arb_pf_replay.py}: \path{e99d81397544d0814b49c8263d61a4a9389dc6688488f4ee0359cc1fc76bfcb2} \\
\emph{Output SHA-256}: \path{s2-selected/sources/numeric/outer-n1536-truepsi.json}: \path{431d567fe714e412da6ba528b02cab0e251cbb616aa1fd8d9c8a11efca67d323} \\
\midrule
\textbf{BV corner, n=2800} \\
$C\in\CertificateInterval{-0.0574908755554524530180814901226778275854309897822443970870}{-0.0570237847462003399142149129742403275854309897822443970868}$ \\
\emph{Source code SHA-256}: \path{s2-selected/sources/numeric/corner_arb.py}: \path{7d5586c43d7951cecd86fdb85fca277efc91e7cf9b85fbd75ada5fb01726d861} \\
\emph{Output SHA-256}: \path{s2-selected/sources/numeric/corner-arb.json}: \path{2bc56db0d72a56afbd88aac7111afc2f2eaa06aa65a4c509a96107765415afc3} \\
\midrule
\textbf{Tail above 12, in D} \\
$[0,0.000047391848857364673265889344422276912]$ \\
\emph{Source code SHA-256}: \path{s2-selected/certify_tail_seed.py}: \path{42565c0188066a45c1f166a49a8fb64c95d6bd43874047bade6af9f856428fef} \\
\emph{Output SHA-256}: \path{s2-selected/TAIL-SEED.json}: \path{0363a93ec43884b70ca0b54e06e71ea8a64210c64910e4e4a3bc3a366df04c53} \\
\midrule
\textbf{Rowwise prefix-repair allowance, in D} \\
$[0,1.154\cdot10^{-54}]$ \\
\emph{Justification}: \\
The preceding tail majorant, within the single $1/20000$ allowance \\
\midrule
\textbf{Joint lower-value assembly} \\
\emph{Outward enclosure assembled from the printed inputs in (A.7)}: \\
$\CertificateInterval{1.6195813569890835255913340433210572922007671408165057072212979}{1.6195813569890835255913340433210572922007671408165057072212983}$ \\
\emph{Full-precision assembler source SHA-256}: \path{s2-selected/verify_interface.py}: \path{fae4b0cb684502ed7d1d75537cb5456868ba83094eb9157f14e30cf58574de71} \\
\emph{Output SHA-256}: \path{s2-selected/INTERFACE-VERIFICATION.json}: \path{839dbd7f3d75d39f52aca3920325d63dd33dcc4346479f0871720d93194fc469} \\
This output uses the retained full-precision prefix endpoint, not the shortened $P_-$ in (A.7); the printed-input enclosure above is evaluated directly from (A.7). \\
\bottomrule
\end{longtable}
}

\paragraph{Independent quadrature.} The calculation recorded in
\path{s2-joint/outputs/ASSEMBLY.json},
SHA-256
\path{21083883403f3903fff3ce781e086a26f931dd809e445578cfcba546295e9e5d},
independently evaluates the PF and corner integrals in the \S~3 functional and gives
\[
\begin{aligned}
 \frac{T_8}{4}&\in
 [1.62063469556829,1.62071322785677]\quad\text{D units},\\
 1.62063469556829-d_{T8}&>0.001053338579.
\end{aligned}\tag{A.8a}
\]
Interval (A.8a) covers the continuous functional and the infinite tail.
The selected $d_{T8}$ uses coarser adverse quadrature endpoints
and also includes the $1/20000$ allowance.
The three improvements are
\[
\begin{aligned}
 \frac{2e^{-\gamma}}4(J_{+,\rm old}-J_{+,\rm new})
   &\simeq0.000944946137955914,\\
 \frac{C_{-,\rm new}-C_{-,\rm old}}4
   &\simeq0.000058392441251611,\\
 \text{tail allowance}&=0.000050000000000000.
\end{aligned}
\]
Here $J_{+,\rm old}$ and $J_{+,\rm new}$ are the printed upper PF
endpoints. Similarly, $C_{-,\rm old}$ and $C_{-,\rm new}$ are the
printed lower corner endpoints. For the new ball we take its outward lower
endpoint. Because of endpoint display rounding, the sum differs from
$1.62063469556829-d_{T8}$ by less than $2\cdot10^{-15}$.
The difference exceeds $0.001053338579$,
so the continuous lower endpoint exceeds $d_{T8}+1/20000$.
The independent prefix-only assembly gives $P_{\rm def}<1.1\cdot10^{-51}$.
The coarser joint prefix-and-tail bound is $10^{-49}$.
The rowwise repair allowance $1.154\cdot10^{-54}$ in (A.7) covers its
own tail-majorant correction and is already inside the single $1/20000$
allowance used in $d_{T8}$. These independent deficits are not added again.
Section~\ref{sec:proof-theorem-1} bounds the finite-$N$ corrections,
including the selector and AA37 terms.
The selected value is (A.8).

The component sources are:

{
\normalsize
\renewcommand{\arraystretch}{1.14}
\begin{longtable}{@{}p{\textwidth}@{}}
\toprule
\textbf{True F/f, 320-bit Arb Taylor models} \\
\emph{Source code SHA-256}: \path{s2-joint/code/ff_models.py}: \path{7a6b137417573ba0e717ffa4444e7fb6671cb1b325ff826cac26630591bd4624} \\
\emph{Output SHA-256/result}: \path{s2-joint/outputs/ff_report.json}: \path{cfb7d637cb623590478e3069b14a2fa2b85aeab1eb296985523b1ccf60f56f01}; $E(s)\le1.302\cdot10^{-55}$ for all $s\ge33$ \\
\midrule
\textbf{PF directed quadrature, 24,000,005 leaves} \\
\emph{Source code SHA-256}: \path{s2-joint/code/jpf.c}: \path{e5a4680d2a6ab4bed93bf8628dae78acabc79bd66bbfcbc6a526ca03cf564973} \\
\emph{Output SHA-256/result}: $J_{PF}\in[0.355795040748992,0.356074784134854]$; four shard output hashes below \\
\midrule
\textbf{BV corner, rigorous nested Arb integration} \\
\emph{Source code SHA-256}: \path{s2-joint/code/corner.py}: \path{7cbc5d276b54b52d657fc52622db1a3542c5079d72a39e3be3d12eec1ddb21f2} \\
\emph{Output SHA-256/result}: \path{s2-joint/outputs/corner.json}: \path{360841b6008cb08be571d23177887df14be286ebde7ba3f3fe2c7571f31be9d8} \\
\midrule
\textbf{Joint directed assembly, including prefix deficit and full tail} \\
\emph{Source code SHA-256}: \path{s2-joint/code/assemble.py}: \path{4b521ae19ef7742e3bcbb4eb28a0f48c2a161195c5ab0e9b4f8456ba6ddf5c45} \\
\emph{Output SHA-256/result}: \path{s2-joint/outputs/ASSEMBLY.json}: \path{21083883403f3903fff3ce781e086a26f931dd809e445578cfcba546295e9e5d}; (A.8a) \\
\bottomrule
\end{longtable}
}

All paths in this table are relative to the supplement root.
The hashes for \path{s2-joint/outputs/jpf_shard0.json} through \path{s2-joint/outputs/jpf_shard3.json} are,
in that order,
\path{14b67605807cfcef9ec96b134fe50b9c8e4b4319c3f33a5ebf99a8eab0cf3c67},
\path{d163c9c8d8d90dbfa5603b49f379867df0912a87e088f0e71924300500bf0fee},
\path{e6f4077dcc2c05317eb2a1b26e4bba519dc7b98cb1960038a2a4885a8bb1f1de},
and
\path{623eebd522d2b99122cf4395cfe760b0d10cc89ec31ac5ca8a70009e8c729b6d}.
The bundle manifest SHA-256 is
\path{d26488e736e5798544e758338acb299d417b74a309568f5ce20f738e399cf96d},
and the manifest contains 16 entries.

The independent assembly gives $P_{\rm def}<1.1\cdot10^{-51}$.
It also bounds the
no-growth tail after $s=33$, and checks the classical A=1 tail argument.
Its PF and corner enclosures are independent of the primary quadrature;
the prefix term $P$ and full volume $V$ are shared inputs.
That section proves the source comparison and the Class-II bounds.
It also bounds the selector, the AA37/G2 terms and the other errors.
The joint four-row endpoint enters once.

\subsection{G1: signed allocation and radial multiplier}

For the exact allocation in \S~4,
\[
 K_G=\frac{2I_{11}+I_{12}+I_{13}}{2e^\gamma},\qquad
 P_r:=-K_GT_r,\qquad \sum_{r=0}^3T_r=\overline G.
                                                               \tag{A.9}
\]
Here $P_r$ is a pure profile target. The corresponding arithmetic
contribution is $Q_{G1-r,0}(N)$; its lower comparison, including the
fixed and vanishing errors, is (4.85). The allocation identity (A.9)
does not assert equality with that arithmetic contribution.
The multiplicities $(2,1,1)$ enter $K_G$ once. The source code is
\path{g1-endpoints/certify_signed_endpoints.py},
with SHA-256 \path{e0eb449e97eedc1ea5aa8ab40de847a1baec1e8316271bf873ead904d845a864}.
Its output \path{g1-endpoints/outputs/result.json}
has SHA-256
\path{0a940b64e5f762fbfc1439709a3092de9749ecfdf448ef467ff52248e1794c69}.
A separate exact rational calculation verifies the retained source intervals.
Its output is
\path{g1-rational/exact-rational-check.json}, with SHA-256
\path{7d7048c800af14d06963bd8fcf4f17ff45517e41fa9fd682bc583ddfb88dd4b6}.
The outward decimal enlargements of these intervals and their product are displayed below.

{
\normalsize
\renewcommand{\arraystretch}{1.14}
\begin{longtable}{@{}p{\textwidth}@{}}
\toprule
\emph{Exact rational outward intervals} \\
\midrule
\textbf{$I_{11}$, including four-prime correction} \\
$\CertificateInterval{0.2208790666923091612026153538231208031879235648368125703408}{0.2208790666923091612026153538231208031879235648368125703424}$ \\
\midrule
\textbf{$I_{12}$} \\
$\CertificateInterval{0.5740619734452483407647685136781950850886922314288474504985}{0.5740619734452483407647685136781950850886922314288474504987}$ \\
\midrule
\textbf{$I_{13}$} \\
$\CertificateInterval{0.4052907329856379464660009714315655397835085202678768274511}{0.4052907329856379464660009714315655397835085202678768274513}$ \\
\midrule
\textbf{$K_G$} \\
$\CertificateInterval{0.3989480791070578466390443009830933753748442947478110112211}{0.3989480791070578466390443009830933753748442947478110112223}$ \\
\midrule
\textbf{$\overline G$} \\
$\CertificateInterval{6.0884403625974206828946573628908904359431953974002641181857}{6.0884403625974206828946573628908904359431953974002641181858}$ \\
\midrule
\textbf{Joint $-K_G\overline G$, in D} \\
$\CertificateInterval{-2.4289715874161197464056136533483979571384269074906618794340}{-2.4289715874161197464056136533483979571384269074906618794266}$ \\
\bottomrule
\end{longtable}
}

These intervals certify G1-0 through G1-3.
All four targets are in D units.
The exact row endpoints below use the narrower exported dyadic $K_G$
and $T_r$ intervals from the retained source.

{
\normalsize
\renewcommand{\arraystretch}{1.14}
\begin{longtable}{@{}p{\textwidth}@{}}
\toprule
\textbf{G1-0}\quad K endpoint for lower bound: upper \\
\emph{Exact lower endpoint}: $\frac{-19564356883289720593180343123839226092697830325553783422353}{98079714615416886934934209737619787751599303819750539264}$ \\
\emph{Exact upper endpoint}: $\frac{-626059420265271058981770979962855234966330570417721069513945}{3138550867693340381917894711603833208051177722232017256448}$ \\
\midrule
\textbf{G1-1}\quad K endpoint for lower bound: lower \\
\emph{Exact lower endpoint}: $\frac{179266388678009569180357614458786189164679660590048624716243}{196159429230833773869868419475239575503198607639501078528}$ \\
\emph{Exact upper endpoint}: $\frac{717065554712038276721430457835144756658718642360194498866525}{784637716923335095479473677900958302012794430558004314112}$ \\
\midrule
\textbf{G1-2}\quad K endpoint for lower bound: upper \\
\emph{Exact lower endpoint}: $\frac{-821300653301490295585240971954986923239672986567876292003797}{392318858461667547739736838950479151006397215279002157056}$ \\
\emph{Exact upper endpoint}: $\frac{-821300653301490295585240971954986923239672986567876292002013}{392318858461667547739736838950479151006397215279002157056}$ \\
\midrule
\textbf{G1-3}\quad K endpoint for lower bound: lower \\
\emph{Exact lower endpoint}: $\frac{270036186059109561485627059871971764166199228966336131972949}{196159429230833773869868419475239575503198607639501078528}$ \\
\emph{Exact upper endpoint}: $\frac{135018093029554780742813529935985882083099614483168065986769}{98079714615416886934934209737619787751599303819750539264}$ \\
\bottomrule
\end{longtable}
}

The exact joint reference value is given below.
{\normalsize
\[
 L_{G1}=-\frac{243950428265194656254206921939947762868871361874409552095097}
 {100433627766186892221372630771322662657637687111424552206336}.
                                                               \tag{A.10}
\]
}
Let $R_{G1}$ be the sum of the four exact row lower endpoints in the
table. The exact difference is
{\normalsize
\[
\begin{gathered}
 L_{G1}-R_{G1}\\
 =\frac{500103}
 {100433627766186892221372630771322662657637687111424552206336}\\
 <4.98\cdot10^{-54}.
\end{gathered}
\]
}
Put
\[
 \delta_{G1}=4.98\cdot10^{-54},\qquad
 E_{G1}=L_{G1}-\delta_{G1}.
\]
The printed decimal intervals are rounded outward separately from the
narrower source intervals used for the exact row certificates.
For an exact check of the endpoint arithmetic, enclose $1/(2e^\gamma)$ by
{\small
\[
\CertificateInterval{0.2807297417834425849120716073954403933828551934625765840770795380225439}{0.2807297417834425849120716073954403933828551934625765840770795380225440}.
\]
}
This rational enclosure comes from the Arb constant evaluation of
the retained source, at 256-bit precision.
We multiply its lower endpoint by $2I_{11,\rm lo}+I_{12,\rm lo}+I_{13,\rm lo}$,
and its upper endpoint by $2I_{11,\rm hi}+I_{12,\rm hi}+I_{13,\rm hi}$.
Exact rational multiplication and outward rounding to $10^{-58}$ then give
the displayed $K_G$ interval.
Its upper endpoint exceeds the upper product by less than
$9.17\cdot10^{-59}$.
The joint interval also contains the displayed interval product.
Put $K_{\rm up}=K_{\rm hi}$ from \eqref{eq:g1-khi}.
Exact multiplication gives
\[
 0<L_{G1}+K_{\rm up}\overline G_{\rm hi}
 <9.36\cdot10^{-58}<\delta_{G1}.
\]
Thus $E_{G1}<-K_{\rm up}\overline G_{\rm hi}$. The exact difference
also gives $E_{G1}<R_{G1}$.
The rational row endpoints keep their independent, sharper certificates.
The selected endpoint is the joint one, $E_{G1}$.
Section~\ref{sec:proof-theorem-1} includes $L_{G1}$ once,
and its error sum includes $\delta_{G1}$ once.
The four row values certify this comparison.
The final sum uses this joint lower endpoint.
The primitive (A.10c) is an earlier alternative.
No completed validated complex quadrature check of its analytic
neighborhoods and branches is used for the selected endpoint.

A separate implementation computes the four-prime correction in two ways
and checks the same $K$ directions. Its source \path{endpoint-recert/code/g1_K.py}
has SHA-256
\path{fd749f9a795f6876c9852c516789a9bbf937ec667c76a06f5662fed480aab52e};
its output \path{endpoint-recert/out/g1_K.json} has SHA-256
\path{e3ea5971a526c195ea33e7eb6b7b0c756d77fa1e0560d3b83eedf1785536cb07}.
The output confirms both $K$ directions. Its residual source \path{endpoint-recert/code/g1_J.py}
has SHA-256
\path{e7e6f75367c3af5d2e97a0c653d866328df962caef8791305203a6ee18be5715};
the output \path{endpoint-recert/out/j_final.json} has SHA-256
\path{7ad65327cda9adf324a5325b664ce3788d6119ea40f31c0c007eb2416f42069c}.
It proves the exact-rational residual enclosure, rounded outward:
\[
 J\in[0.964498,0.967120],\qquad J<0.968435989360.                                \tag{A.10a}
\]
The residual-profile $\tau$-effect is at most
$5\cdot10^4\tau\le5\cdot10^{-86}$. The selected residual slack absorbs it.
The comparisons give all four T allocations and lower targets,
with signs $(-,+,-,+)$.
The $K$ directions are
$(\mathrm{upper},\mathrm{lower},\mathrm{upper},\mathrm{lower})$,
with the weights $(2,1,1)$ entering once. Section~4 proves the source application for
$C_b=334$ and its error bounds.

\paragraph{The 4096-cell upper-corner calculation.}
This paragraph gives the 4096-cell trapezoidal calculation of $\overline G$ used in (4.79).
The residual enclosure (A.10a) comes from the separate Taylor-model source and output identified above.
The true delay functions are

\[
 F(s)={2e^\gamma\over s}\ (1\le s\le3),\quad
 f(s)=0\ (s\le2),\quad
 f(s)={2e^\gamma\log(s-1)\over s}\ (2\le s\le3),
\]

and thereafter $(sF)'=f(s-1)$, $(sf)'=F(s-1)$.
On open integer intervals, their second derivatives are bounded by 12.
The quadrature uses trapezoidal panels of width $2^{-16}$ aligned at integers,
with integral remainder at most $h^3$.
Arb carries all propagated radii, and outward rational rounding uses the
$2^{-128}$ lattice.
Monotonicity encloses $\Phi(a)=\int_1^a(1-f(s))\,ds$. Beyond 30, its tail is bounded by
$(a-30)(1-f(30))$; no constant table extrapolation is used.

For fixed rational $t_1,t_2$, the $t_3$ range is split at each kink,
activation switch and pairwise affine crossing in (4.70).
On each segment $a_3=c/t_3+k$, and

\[
 \int_p^q{1-f(c/t+k)\over t^2}\,dt
     ={\Phi(c/p+k)-\Phi(c/q+k)\over c}\quad(c\ne0).
 \tag{A.10b}\label{eq:g1-phi-primitive}
\]

The constant-argument case is direct. Each segment integral lies
between 0 and $1/p-1/q$. The residual
$\Theta_*-t_1-t_2-t_3$ is nonincreasing in $t_1,t_2$.
Its affine derivatives are negative, and the $\psi$ active sets are down-sets.
Hence the inner integral is nondecreasing in both outer variables.
The outer integration uses $t=ue^x$, outward-enclosed logarithmic cells, and upper
and lower corner sums. Diagonal triangles have half the squared cell width.
On a grid of 4096 cells, the outward-rounded sums are accumulated as exact
integers. This proves the directed enclosure (4.79).

For an independent evaluation of the four-prime correction, put
$z=v/(1-t)$. A primitive of
$\log((1-2z)/z)/(z(1-z))$ is

\[
\begin{aligned}
 P(z)={}&-\operatorname{Li}_2(2z)-\tfrac12\log^2z
 -\tfrac12\log^2(1-z)-\log2\log(1-z)\\
 &-\operatorname{Li}_2(1/[2(1-z)])
       +\log z\log(1-z)+\operatorname{Li}_2(z).
\end{aligned}
 \tag{A.10c}\label{eq:g1-dilog-primitive}
\]

The formula is checked by differentiation. The analytic neighborhoods
and the real branches would require a validated complex-quadrature check
before this primitive could serve as an endpoint certificate.
Nested integration of the original double integral gives an
independent enclosure. At 200-bit precision, the signed enclosures are:

{
\normalsize
\setlength{\tabcolsep}{3pt}
\begin{longtable}{@{}p{\textwidth}@{}}
\toprule
\textbf{$I_{11}$} \\
\emph{Abbreviated centre}: $0.22087906669230916120261535382312080318$ \\
\emph{Width less than}: $1.476\cdot10^{-57}$ \\
\midrule
\textbf{$I_{12}$} \\
\emph{Abbreviated centre}: $0.57406197344524834076476851367819508508$ \\
\emph{Width less than}: $6.348\cdot10^{-59}$ \\
\midrule
\textbf{$I_{13}$} \\
\emph{Abbreviated centre}: $0.40529073298563794646600097143156553978$ \\
\emph{Width less than}: $3.672\cdot10^{-59}$ \\
\midrule
\textbf{$K_G$} \\
\emph{Abbreviated centre}: $0.39894807910705784663904430098309337537$ \\
\emph{Width less than}: $8.588\cdot10^{-58}$ \\

\bottomrule
\end{longtable}
}


These centres are not directed endpoints. The source-derived widths above
refer to the narrower retained intervals. For the printed decimal
integral intervals, the exact outward upper endpoint is

\[
 K_{\rm hi}=
 0.3989480791070578466390443009830933753748442947478110112223.
 \tag{A.10d}\label{eq:g1-khi}
\]


{
\normalsize
\setlength{\tabcolsep}{3pt}
\begin{longtable}{@{}p{\textwidth}@{}}
\toprule
\textbf{G1-0} \\
\emph{Abbreviated centre}: $-199.47403966257515881115583710023774090$ \\
\emph{$K_G$ direction}: upper \\
\midrule
\textbf{G1-1} \\
\emph{Abbreviated centre}: $+913.88106797076246277430172018777105602$ \\
\emph{$K_G$ direction}: lower \\
\midrule
\textbf{G1-2} \\
\emph{Abbreviated centre}: $-2093.45188381184444127509736016826295322$ \\
\emph{$K_G$ direction}: upper \\
\midrule
\textbf{G1-3} \\
\emph{Abbreviated centre}: $+1376.61588391624101756554586342738124014$ \\
\emph{$K_G$ direction}: lower \\

\bottomrule
\end{longtable}
}



The narrower source-derived two-sided joint width is below $5.248\cdot10^{-57}$.
The exact gap is the fraction displayed after (A.10).

\subsection{Switched mass, BV endpoints and finite threshold}

Write $H_L:=I_{HL}$ as in (5.32), equivalently (1.12), with the parameters
and cell weight (5.30)--(5.31). The measure is the coordinate measure $dt_1\cdots dt_{k-1}$,
where $t_k=1-\sum_{i<k}t_i$.
It carries no factorial or Euclidean surface-measure factor.
Then

\[
 0\le H_L<0.0778506,\qquad -2H_L>-0.1557012.            \tag{A.11}
\]
Both implementations cover every positive-weight stratum $5\le k\le11$,
and the coverage is exact rational.
Their upper endpoints lie in
\[
 [0.0778505820459738841964296666351,
   0.0778505820459738841964296666353].                   \tag{A.12}
\]
The selected source is
\path{switched-mass/original/scripts/refine.py}, with SHA-256
\path{83ff0a33baf95be8400a69c466965baae130b3cd788acd2d83b83fd2cc45398b}.
The separate calculation encloses the same switched mass.
The independent source
\path{switched-mass/independent/independent_replay.py}
has SHA-256
\path{c275b1904b506cf82315d4e27f44f1949487e67044923501fd0628738a7744ad}.
Their common compact output
\path{switched-mass/replay_summary.json}
has SHA-256
\path{f0dac63f31707d0ecfe9ffd08fdfe21d4e5b2718b19a83af713ba518d2077e06}.
Exact rational coverage and overlap identities give the scalar input for
the switched term of \S~5.

For the S6 BV row, put
\[
 J(\kappa)=\int_{\tau_S}^{1/3}\frac{dt}{t(1/2-\kappa-t)}.
\]
Exact logarithms and independent Arb arithmetic give the enclosures.
{\normalsize
\[
\begin{gathered}
 J(0)\in\\\CertificateInterval{1.3148581959157321533513247276782285310758770468886237619584}{1.3148581959157321533513247276782285310758770468886237619585},\\
 J(10^{-6})\in\\\CertificateInterval{1.3148646802099601102520747017315401230927649569798117109688}{1.3148646802099601102520747017315401230927649569798117109689},\\
 -[J(10^{-6})-J(0)]\in\\
 \CertificateInterval{-6.48429422795690074997405331159202\cdot10^{-6}}{
 -6.48429422795690074997405331159201\cdot10^{-6}}.
 \end{gathered}\tag{A.13}
\]
}
Here $S_6=4J$, so the nominal D contribution is $-J(0)$.
The increment in (A.13) enters once at level $1/2-10^{-6}$ of
\S~5. The source
\path{row-enclosures/cert_s4_s7.py}
has SHA-256
\path{d98866b45f1ddb8f3f93411bfe09cb11c62569cc80fa775a77a3eab1ad4384aa};
its \path{row-enclosures/result.json} has SHA-256
\path{0eb5fa52ed180a071633e50f0d099f5c4e1f9bb001add71765fe1c2f033b8f7a}.
The independent 256-bit calculation is
\path{bv-increment/checks/charge_check.py},
with SHA-256
\path{1a7be201efe55e801c2ea7c2dc0807ac543d323838165f259b982a064e460ac6},
and its output \path{bv-increment/checks/charge_check.out}
has SHA-256
\path{e99ef910bb85c1bd2119314c37a0deddcf1bc11f43c1d4fe50e8aaf956748015}.
It verifies the exact-log enclosures in (A.13).

\paragraph{S5 BV-top normalization.} The moved-cell value is
enclosed by
\[
 B(0.03658166423843258896399704402764682632822246022922592990,
   8.84\cdot10^{-57}).                                 \tag{A.14}
\]
This raw row is $4J_5(0)$, with master coefficient $-1$ by
(5.15)--(5.20a). Its D contribution is $-J_5(0)$, with exact interval
{\normalsize
\[
\begin{gathered}
 -J_5(0)\in\\
-\frac14 B\left(\begin{gathered}0.03658166423843258896399704402764682632822246022922592990,\\
             8.84\cdot10^{-57}\end{gathered}\right),
\end{gathered}\tag{A.14a}
\]
}
which is about $-0.00914541605960815$ D. Section~5 proves the
normalization. S5-low uses its direct-cell endpoint, and S5 BV-top uses
the moved-cell exact-log value (A.14).
Their supports are disjoint at $t_*$,
and each contribution enters once.

For comparison, if S5-top used $\kappa=10^{-6}$, its negative coefficient
would require an upper bound for $J_5(\kappa)$. Since $t\ge t_*$ and
$1/2-t\ge1/2-\tau_S$,
\[
\begin{aligned}
 0\le J_5(\kappa)-J_5(0)
 &=\int_{t_*}^{\tau_S}
   \frac{\kappa\,dt}{t(1/2-t)(1/2-t-\kappa)}\\
 &\le\frac{(\tau_S-t_*)\kappa}
 {t_*(1/2-\tau_S)(1/2-\tau_S-\kappa)}\\
 &=\frac{25088}{672019919285}<3.734\cdot10^{-8}<4\cdot10^{-8}.
\end{aligned}
\]
Thus $-J_5(\kappa)\ge-J_5(0)-E_{5,\kappa}$,
with $E_{5,\kappa}=4\cdot10^{-8}$ in D units.
Section~\ref{sec:proof-theorem-1} subtracts this error once.
The separate conservative BV allowance is
$g_{\rm BV}=5.75\cdot10^{-6}=E_{5,\kappa}+g_{R1}
=4\cdot10^{-8}+5.71\cdot10^{-6}$.
The R1 increment itself is already inside the selected S2 endpoint.
Equations (A.14)--(A.14a) give the nominal $\kappa=0$ value.

For an independent exact-rational S5-top certificate, the logarithmic
ratio in (5.20) is
$R=392787/390995$, so $z=(R-1)/(R+1)=896/391891$. Define
\[
 P_5=16\sum_{k=0}^{15}\frac{z^{2k+1}}{2k+1},\qquad
 R_5=\frac{16z^{33}}{33(1-z^2)}.
\]
Then $P_5<4J_5(0)<P_5+R_5$ by the positive atanh series.
Exact rational comparisons give $R_5<3.447\cdot10^{-88}$ and
{\normalsize
\[
\begin{gathered}
 [P_5,P_5+R_5]\subset\\
B(0.03658166423843258896399704402764682632822246022922592990,
   8.84\cdot10^{-57}).
\end{gathered}
\]
}
This proves the directed $-1/4$ normalization. The printed rational
series is the certificate. Section~6.3 proves its source application
and arithmetic error bounds.

The finite-epsilon target is
\[
\begin{aligned}
 T&=8\log\frac{50009}{39991},\\
 1.78838848345960807361670991332&<T
 <1.78838848345960807361670991333.
\end{aligned}\tag{A.15}
\]
It is proved by exact rational arithmetic. For
$z=5009/45000$, set
\[
 P=16\sum_{k=0}^{23}\frac{z^{2k+1}}{2k+1},\qquad
 R=\frac{16z^{49}}{49(1-z^2)}.
\]
Then $P<T<P+R$, and exact rational comparisons give (A.15).
This series and its tail bound certify (A.15),
and its upper endpoint enters the final strict comparison.

\subsection{Selected numerical endpoints}

The table below lists the selected values.
The full endpoints and proofs are in A.2--A.6.
\begin{longtable}{@{}p{9em}p{\dimexpr\textwidth-9em-4\tabcolsep\relax}@{}}
\toprule
Contribution & Selected value and direction \\
\midrule
S1 & Joint lower endpoint (A.4), with multiplier $3/4$ \\
S2 & Joint lower value (A.8), including the $1/20000$ tail \\
S3/S8/S9 & Absolute Option-A endpoints; signs $-,+,+$ \\
S4/S5-low/S7 & Direct-row upper endpoints; multiplier $-1/4$ \\
S10-Ret & Retained-row lower endpoint; multiplier $1/4$ \\
G1 & $E_{G1}=L_{G1}-\delta_{G1}$, once \\
S5-top & Raw-to-D factor $-1/4$; separate error $E_{5,\kappa}$ \\
S6 & Nominal value $-J(0)$ and increment (A.13) \\
Switched term & $-2H_L>-0.1557012$ by (A.11) \\
Threshold & Upper endpoint (A.15) \\
\bottomrule
\end{longtable}
Sections~2--5 prove the source applications,
and Section~6.3 combines the terms with the signs of (6.1).
}
